\Needspace{6\baselineskip}
\begin{center}\hypertarget{native-251-article179}{}\pdfbookmark[2]{179.}{article-179}179.\end{center}
Wir entwickeln nun die Folgerungen aus diesem Gesetz. Damit das Produkt $\Omega=h^\mu\pi^{-\frac{1}{2}\mu}e^{-hh(vv+v'v'+v''v''+\mathord{.}\,\mathord{.})}$ ein Maximum wird, muss offenbar die Summe $vv+v'v'+v''v''+\mathord{.}\,\mathord{.}$ ein Minimum werden. \emph{Das wahrscheinlichste System von Werten für die Unbekannten $p$, $q$, $r$, $s$ usw. ist demnach dasjenige, bei dem die Quadrate der Differenzen zwischen den beobachteten und den berechneten Werten der Funktionen $V$, $V'$, $V''$ usw. die kleinste Summe ergeben,} sofern bei allen Beobachtungen derselbe Genauigkeitsgrad vorauszusetzen ist.

Dieses Prinzip findet bei allen Anwendungen der Mathematik auf die Naturphilosophie sehr häufig Verwendung. Es soll überall mit demselben Recht als Axiom gelten, mit dem das arithmetische Mittel aus mehreren beobachteten Werten ein und derselben Größe als ihr wahrscheinlichster Wert angenommen wird.

Das Prinzip lässt sich nun ohne Schwierigkeit auf Beobachtungen von \emph{ungleicher} Genauigkeit ausdehnen. Das Maß der Genauigkeit der Beobachtungen, aus denen $V=M$, $V'=M'$, $V''=M''$ usw. gefunden wurden, werde der Reihe nach durch $h$, $h'$, $h''$ usw. ausgedrückt. Man nehme also an, dass Fehler, die diesen Größen umgekehrt proportional sind, bei diesen Beobachtungen ebenso leicht auftreten konnten. Das ist offenbar gleichbedeutend damit, dass durch Beobachtungen gleicher Genauigkeit (mit dem Maß $=1$) für die Funktionen $hV$, $h'V'$, $h''V''$ usw. unmittelbar die Werte $=hM$, $h'M'$, $h''M''$ usw. gefunden worden wären. \SourceLeaf{252}Das wahrscheinlichste System von Werten für die Größen $p$, $q$, $r$, $s$ usw. ist daher dasjenige, bei dem die Summe $hhvv+h'h'v'v'+h''h''v''v''+\text{usw.}$, also \emph{die Summe der Quadrate der Differenzen zwischen den tatsächlich beobachteten und den berechneten Werten, wobei jede Differenz mit der Zahl multipliziert wird, die den Grad der Genauigkeit misst, ein Minimum wird.} Dabei brauchen sich die Funktionen $V$, $V'$, $V''$ usw. nicht einmal auf gleichartige Größen zu beziehen. Sie können auch verschiedenartige Größen darstellen (zum Beispiel Bogensekunden und Zeitsekunden), sofern sich das Verhältnis der Fehler abschätzen lässt, die bei den einzelnen Beobachtungen ebenso leicht auftreten konnten.
